\section{Алгоритмические языки. Представление о языках программирования высокого уровня. Пакеты прикладных программ}

\subsection{Алгоритм и алгоритмический язык}

\textbf{Алгоритм} --- конечная последовательность однозначно определённых действий, приводящая от исходных данных к результату за конечное число шагов.

Обычно выделяют следующие основные свойства алгоритма:
\begin{itemize}
	\item \textbf{дискретность} --- алгоритм состоит из отдельных шагов;
	\item \textbf{определённость} --- каждое действие должно быть однозначно задано;
	\item \textbf{конечность} --- для допустимых исходных данных выполнение должно завершаться за конечное число шагов;
	\item \textbf{результативность} --- после завершения должен быть получен некоторый результат;
	\item \textbf{массовость} --- алгоритм предназначен не только для одного набора данных, а для некоторого класса однотипных задач.
\end{itemize}

Алгоритм можно задавать различными способами:
\begin{itemize}
	\item словесным описанием;
	\item математическими формулами;
	\item блок-схемой;
	\item псевдокодом;
	\item программой на формальном языке.
\end{itemize}

\textbf{Алгоритмический язык} --- формальный язык, предназначенный для записи алгоритмов в форме, допускающей их однозначное понимание и, как правило, автоматическое выполнение или преобразование.

Алгоритмический язык задаёт:
\begin{itemize}
	\item алфавит допустимых символов;
	\item правила образования конструкций языка;
	\item способы представления данных;
	\item управляющие конструкции;
	\item правила выполнения записанных конструкций.
\end{itemize}

Исторически понятие алгоритмического языка является несколько более широким, чем понятие языка программирования. Алгоритм можно записывать, например, на формализованном псевдокоде, который не предназначен непосредственно для выполнения компьютером. \textbf{Язык программирования} представляет собой формальный алгоритмический язык, для которого определены средства трансляции и выполнения программ на вычислительной системе.

Типичная программа содержит три основных составляющих:
\begin{enumerate}
	\item описание данных;
	\item описание операций над ними;
	\item описание порядка выполнения этих операций.
\end{enumerate}

Для структурного программирования базовыми управляющими конструкциями являются:
\begin{itemize}
	\item \textbf{следование} --- последовательное выполнение операторов;
	\item \textbf{ветвление} --- выбор одного из вариантов вычисления;
	\item \textbf{цикл} --- многократное выполнение некоторого набора действий.
\end{itemize}


\subsection{Синтаксис и семантика языка}

Описание любого формального языка включает две принципиально различные части: \textbf{синтаксис} и \textbf{семантику}.

\textbf{Синтаксис} определяет, какие последовательности символов являются правильными конструкциями языка.

К синтаксису относятся:
\begin{itemize}
	\item лексемы языка: идентификаторы, ключевые слова, литералы, знаки операций;
	\item правила записи выражений;
	\item правила объявления переменных и функций;
	\item правила построения операторов;
	\item структура программных модулей.
\end{itemize}

Например, выражение

\[
a = b + c;
\]

может быть синтаксически корректным в одном языке, тогда как отсутствие точки с запятой в том же месте может приводить к синтаксической ошибке.

Формальное описание синтаксиса часто строится при помощи грамматик, например формы Бэкуса--Наура.

\textbf{Семантика} определяет смысл синтаксически корректной конструкции, то есть результат её выполнения.

Например, выражение

\[
x = x + 1
\]

с точки зрения математики является противоречивым равенством, но в императивном языке программирования оно означает присваивание переменной $x$ значения, на единицу большего предыдущего.

Различают:
\begin{itemize}
	\item \textbf{статическую семантику} --- свойства программы, проверяемые до её выполнения;
	\item \textbf{динамическую семантику} --- поведение программы непосредственно во время выполнения.
\end{itemize}

К статическим проверкам могут относиться:
\begin{itemize}
	\item соответствие типов;
	\item область видимости идентификаторов;
	\item корректность количества и типов аргументов функции.
\end{itemize}

Важным элементом языка является \textbf{система типов}. Тип определяет множество допустимых значений и операций над ними.

Распространённые типы данных:
\begin{itemize}
	\item целые числа;
	\item вещественные числа;
	\item логические значения;
	\item символы и строки;
	\item массивы;
	\item записи и структуры;
	\item указатели и ссылки;
	\item пользовательские типы и классы.
\end{itemize}

Системы типов могут быть:
\begin{itemize}
	\item \textbf{статическими}, когда большинство типов известно до выполнения программы;
	\item \textbf{динамическими}, когда тип значения определяется непосредственно во время выполнения.
\end{itemize}


\subsection{Языки низкого и высокого уровня}

Уровень языка характеризует степень абстрагирования от конкретной архитектуры вычислительной машины.

\textbf{Языки низкого уровня} тесно связаны с системой команд процессора.

К ним относят:
\begin{itemize}
	\item машинный код;
	\item языки ассемблера.
\end{itemize}

Машинная программа непосредственно представляется последовательностью команд процессора.

Ассемблер использует символические обозначения машинных команд, например:

\begin{verbatim}
	MOV AX, BX
	ADD AX, 1
\end{verbatim}

Преимущества низкоуровневого программирования:
\begin{itemize}
	\item непосредственный контроль над аппаратурой;
	\item возможность точного управления памятью;
	\item потенциально высокая производительность;
	\item возможность использования специфических команд процессора.
\end{itemize}

Недостатки:
\begin{itemize}
	\item сильная зависимость программы от архитектуры;
	\item высокая трудоёмкость разработки;
	\item сложность сопровождения;
	\item повышенная вероятность ошибок.
\end{itemize}

\textbf{Языки высокого уровня} предоставляют программисту абстракции, не связанные непосредственно с машинными командами.

К ним относятся, например:
\texttt{Fortran}, \texttt{C++}, \texttt{Java}, \texttt{Python}, \texttt{C\#}.

Основные особенности высокоуровневых языков:
\begin{itemize}
	\item наличие развитых типов данных;
	\item возможность использования функций и модулей;
	\item наличие структур управления;
	\item автоматизация части работы с памятью;
	\item библиотеки стандартных алгоритмов и структур данных;
	\item относительная переносимость исходного кода между различными платформами.
\end{itemize}

Высокоуровневый язык позволяет описывать задачу в терминах, более близких к предметной области, а не к системе команд процессора.

Следует отметить, что деление на уровни условно. Например, \texttt{C} и \texttt{C++} являются языками высокого уровня, но предоставляют достаточно низкоуровневые средства работы с памятью и аппаратурой.


\subsection{Парадигмы программирования}

\textbf{Парадигма программирования} --- совокупность принципов и способов организации вычислений и структуры программы.

Современные языки часто являются \textbf{мультипарадигменными}, то есть поддерживают одновременно несколько подходов.

\paragraph{Императивное программирование}

Программа рассматривается как последовательность команд, изменяющих состояние вычислительной системы.

Основной механизм --- присваивание:

\[
x \leftarrow x + 1.
\]

К императивным языкам относятся, например, \texttt{C}, \texttt{Fortran}, \texttt{Pascal}.

\paragraph{Процедурное программирование}

Является развитием императивного подхода. Программа разбивается на процедуры и функции.

Основные преимущества:
\begin{itemize}
	\item декомпозиция задачи;
	\item повторное использование кода;
	\item локализация данных и алгоритмов.
\end{itemize}

\paragraph{Объектно-ориентированное программирование}

Программа строится как система взаимодействующих объектов.

Основными понятиями являются:
\begin{itemize}
	\item класс;
	\item объект;
	\item инкапсуляция;
	\item наследование;
	\item полиморфизм.
\end{itemize}

Объект объединяет данные и операции над ними.

Например, объект, моделирующий матрицу, может хранить её элементы и предоставлять методы умножения, вычисления нормы или решения системы линейных уравнений.

\paragraph{Функциональное программирование}

Вычисление рассматривается как вычисление значений функций.

Характерные особенности:
\begin{itemize}
	\item широкое использование функций высших порядков;
	\item стремление уменьшить количество изменяемых состояний;
	\item использование рекурсии;
	\item удобство композиции функций.
\end{itemize}

\paragraph{Логическое программирование}

Программа описывает набор фактов и логических правил, а вычисление заключается в логическом выводе.

Классическим примером является язык \texttt{Prolog}.

\paragraph{Декларативный подход}

В декларативных языках программист преимущественно указывает, \textit{что} требуется получить, не задавая детально, \textit{как} именно выполнить вычисление.

К декларативным языкам предметных областей можно отнести, например, язык запросов \texttt{SQL}.


\subsection{Трансляторы: компиляторы и интерпретаторы}

Для выполнения программы, записанной на языке высокого уровня, необходимо преобразовать её в форму, понятную вычислительной системе.

\textbf{Транслятор} --- программа, преобразующая текст с одного формального языка на другой.

Основные виды трансляторов:
\begin{itemize}
	\item ассемблер;
	\item компилятор;
	\item интерпретатор.
\end{itemize}

\textbf{Ассемблер} преобразует программу на языке ассемблера в машинный код.

\textbf{Компилятор} преобразует исходную программу целиком или крупными единицами в машинный код либо в некоторое промежуточное представление.

Типичная схема компиляции:

\[
\text{исходный текст}
\rightarrow
\text{лексический анализ}
\rightarrow
\text{синтаксический анализ}
\rightarrow
\text{семантический анализ}
\rightarrow
\text{оптимизация}
\rightarrow
\text{генерация кода}.
\]

На этапе \textbf{лексического анализа} последовательность символов разбивается на лексемы.

На этапе \textbf{синтаксического анализа} проверяется соответствие программы грамматике языка и строится её синтаксическое представление.

При \textbf{семантическом анализе} проверяются типы, объявления идентификаторов и другие смысловые ограничения.

После этого компилятор может выполнять различные оптимизации и генерировать целевой код.

Отдельным этапом является \textbf{компоновка} (\textit{linking}), при которой объектные модули программы и необходимые библиотеки объединяются в исполняемый файл.

\textbf{Интерпретатор} непосредственно выполняет программу или её промежуточное представление, не создавая заранее самостоятельный машинный исполняемый файл для всей программы.

Схематично:

\[
\text{исходная программа}
\rightarrow
\text{интерпретатор}
\rightarrow
\text{выполнение}.
\]

На практике граница между компиляцией и интерпретацией не является строгой.

Например, распространена схема

\[
\text{исходный код}
\rightarrow
\text{байт-код}
\rightarrow
\text{виртуальная машина}
\rightarrow
\text{машинный код}.
\]

Также используется \textbf{JIT-компиляция} (\textit{Just-In-Time}), при которой отдельные участки программы переводятся в машинный код непосредственно во время выполнения.

Поэтому язык сам по себе не является строго ``компилируемым'' или ``интерпретируемым'': это прежде всего свойство конкретной реализации языка.


\subsection{Среда выполнения и библиотеки}

Работа программы определяется не только самим языком и его транслятором. Важную роль играет \textbf{среда выполнения} (\textit{runtime environment}).

Среда выполнения может обеспечивать:
\begin{itemize}
	\item загрузку программы;
	\item управление памятью;
	\item организацию вызовов функций;
	\item обработку исключений;
	\item ввод и вывод;
	\item взаимодействие с операционной системой;
	\item динамическую загрузку библиотек;
	\item выполнение промежуточного кода;
	\item многопоточность.
\end{itemize}

Память выполняющейся программы условно можно разделить на несколько областей:
\begin{itemize}
	\item область машинного кода;
	\item область глобальных и статических данных;
	\item \textbf{стек} вызовов функций;
	\item \textbf{кучу} для динамически создаваемых объектов.
\end{itemize}

В одних языках управление динамической памятью выполняется программистом, в других значительная часть этой работы возлагается на среду выполнения и \textbf{сборщик мусора}.

\textbf{Библиотека} --- набор готовых программных компонентов, предназначенных для повторного использования.

Стандартные библиотеки языка обычно предоставляют:
\begin{itemize}
	\item контейнеры и структуры данных;
	\item математические функции;
	\item работу со строками;
	\item файловый ввод-вывод;
	\item средства многопоточности;
	\item сетевые функции.
\end{itemize}

Кроме стандартной библиотеки существуют специализированные библиотеки.

В научных вычислениях распространены библиотеки для:
\begin{itemize}
	\item линейной алгебры;
	\item решения ОДУ и УЧП;
	\item оптимизации;
	\item статистического анализа;
	\item параллельных вычислений;
	\item визуализации результатов.
\end{itemize}

Использование библиотек позволяет не реализовывать повторно стандартные алгоритмы и применять хорошо протестированные и оптимизированные программные компоненты.


\subsection{Пакеты прикладных программ}

\textbf{Пакет прикладных программ (ППП)} --- комплекс взаимосвязанных программных средств, предназначенных для решения некоторого класса прикладных задач.

В отличие от отдельной библиотеки, ППП обычно предоставляет пользователю более высокоуровневую законченную среду решения задач.

Пакет может включать:
\begin{itemize}
	\item пользовательский интерфейс;
	\item язык команд или встроенный язык программирования;
	\item набор алгоритмов;
	\item библиотеки;
	\item средства хранения данных;
	\item средства визуализации;
	\item средства импорта и экспорта;
	\item документацию.
\end{itemize}

ППП можно классифицировать по назначению.

\paragraph{Пакеты общего назначения}

Используются для типовых пользовательских задач:
\begin{itemize}
	\item текстовые процессоры;
	\item электронные таблицы;
	\item системы управления базами данных.
\end{itemize}

\paragraph{Математические пакеты}

Предназначены для аналитических и численных вычислений.

Можно выделить:
\begin{itemize}
	\item системы компьютерной алгебры;
	\item системы численных вычислений;
	\item статистические пакеты;
	\item средства оптимизации.
\end{itemize}

Например, математические системы могут выполнять:
\begin{itemize}
	\item операции линейной алгебры;
	\item решение систем уравнений;
	\item численное интегрирование;
	\item решение дифференциальных уравнений;
	\item оптимизацию;
	\item обработку экспериментальных данных;
	\item построение графиков.
\end{itemize}

\paragraph{Пакеты математического моделирования}

Для сложного моделирования обычно используется типовая технологическая цепочка:

\[
\text{геометрия}
\rightarrow
\text{математическая модель}
\rightarrow
\text{дискретизация}
\rightarrow
\text{решатель}
\rightarrow
\text{постобработка}.
\]

Такие комплексы могут содержать:
\begin{itemize}
	\item генераторы расчётных сеток;
	\item конечно-разностные или конечно-элементные решатели;
	\item решатели систем линейных и нелинейных уравнений;
	\item средства параллельных вычислений;
	\item средства визуализации полей и сеток.
\end{itemize}

Примерами классов специализированных ППП являются:
\begin{itemize}
	\item системы конечно-элементного анализа;
	\item пакеты вычислительной гидродинамики;
	\item системы компьютерной алгебры;
	\item статистические системы;
	\item системы автоматизированного проектирования.
\end{itemize}

Важным преимуществом ППП является возможность решать сложные прикладные задачи без реализации всей вычислительной инфраструктуры с нуля.

Однако применение готового пакета не освобождает исследователя от необходимости понимать:
\begin{itemize}
	\item используемую математическую модель;
	\item численный метод;
	\item ограничения решателя;
	\item погрешности вычислений;
	\item условия применимости полученного результата.
\end{itemize}

Иначе возникает ситуация, когда программный комплекс выдаёт численный результат, но пользователь не способен оценить его корректность.


\subsection{Критерии выбора языка и программного пакета}

Выбор языка программирования определяется не существованием некоторого ``лучшего языка'', а требованиями конкретной задачи.

Основные критерии выбора языка:
\begin{enumerate}
	\item \textbf{соответствие предметной области};
	\item \textbf{производительность};
	\item \textbf{наличие необходимых библиотек};
	\item \textbf{простота разработки и сопровождения};
	\item \textbf{переносимость};
	\item \textbf{поддержка параллельных вычислений};
	\item \textbf{взаимодействие с другими языками и программами};
	\item \textbf{качество средств тестирования и отладки};
	\item \textbf{наличие специалистов и развитость экосистемы}.
\end{enumerate}

Например, при численном моделировании может использоваться комбинация нескольких языков:
\begin{itemize}
	\item высокопроизводительное вычислительное ядро реализуется на компилируемом языке;
	\item подготовка данных и автоматизация расчётов выполняются на скриптовом языке;
	\item специализированные библиотеки используются для линейной алгебры, сеток и параллельных вычислений.
\end{itemize}

При выборе пакета прикладных программ следует учитывать:
\begin{enumerate}
	\item наличие необходимых математических моделей;
	\item реализованные численные методы;
	\item возможность задания требуемых начальных и граничных условий;
	\item устойчивость и точность решателей;
	\item возможности построения и импорта сеток;
	\item масштабируемость на многопроцессорных системах;
	\item средства визуализации и постобработки;
	\item возможность автоматизации серии расчётов;
	\item поддерживаемые форматы данных;
	\item возможность интеграции с собственным программным кодом;
	\item наличие документации и технической поддержки;
	\item стоимость и лицензионные ограничения;
	\item воспроизводимость вычислительных результатов.
\end{enumerate}

Таким образом, язык программирования и пакет прикладных программ решают разные уровни задачи. Язык представляет средства построения алгоритмов и программ, библиотеки предоставляют готовые программные компоненты, а пакет прикладных программ является более законченной системой решения определённого класса прикладных задач.

Для математического моделирования типична следующая иерархия:

\[
\boxed{
	\text{математическая модель}
	\rightarrow
	\text{алгоритм}
	\rightarrow
	\text{программная реализация}
	\rightarrow
	\text{вычислительный результат}
}
\]

При этом корректность конечного результата определяется не только качеством программного обеспечения, но и адекватностью модели, корректностью алгоритма и правильностью выбора численных методов.

\newpage